\chapter{Panic Disorder Test}
\section*{What Is This Test?} Panic-disorder assessment evaluates recurrent unexpected panic attacks and persistent concern or avoidance.\section*{Alternative Names} Panic attack assessment; panic screening; panic symptom questionnaire.\section*{Principle} Clinical interview establishes attack features, recurrence, impairment, and alternative medical or psychiatric causes.\section*{Why This Test Is Ordered} It is used for sudden episodes of intense fear with palpitations, breathlessness, chest discomfort, dizziness, or fear of dying.\section*{Primary vs Secondary Test} Clinical assessment is primary; ECG, glucose, thyroid, toxicology, or other tests are targeted to symptoms and risk.\section*{How This Test Is Performed} The clinician reviews episodes, triggers, substances, medicines, trauma, mood, and safety.\section*{What Preparation Is Needed} None; provide timing and symptom details. New severe chest pain, fainting, or breathing difficulty needs urgent assessment.\section*{Understanding Results} Panic symptoms can resemble cardiopulmonary or neurologic disease; diagnosis requires exclusion of likely dangerous alternatives.\section*{Follow-up Tests} ECG, laboratory tests, psychiatric assessment, and screening for depression or substance use may follow.\section*{References} \begin{enumerate}\item MedlinePlus. Panic Disorder Test.\item American Psychiatric Association. Panic disorder diagnostic guidance.\end{enumerate}\clearpage\section*{Understanding Results and Follow-up}
The laboratory report should identify the specimen, method, units, reference interval or decision limit, and whether the result is positive, negative, abnormal, or inconclusive. Reference intervals vary with age, sex, pregnancy, specimen type, assay, and laboratory; a result outside a range is not automatically a diagnosis, and a result within a range does not always exclude disease. Discuss unexpected results with the ordering clinician, who can decide whether repeat or confirmatory testing, treatment, or referral is appropriate. Do not change medicines or delay urgent care based on an isolated result.
\section*{Regulatory Status and Authorization Date}This chapter describes a test category, not a specific commercial product. FDA, CDSCO, EU IVDR, and MHRA approval or registration dates therefore cannot be assigned without the assay manufacturer, product name, instrument, and intended use. For laboratory-developed tests, an individual agency approval date may not exist. See the Preface for the interpretation of regulatory status.
\clearpage
